\documentclass{article}
\usepackage{booktabs,threeparttable,tabularx}
\begin{document}

\begin{table}[htbp]
\centering
\caption{Panel Unit Root Test: Levin--Lin--Chu (2002)}
\label{tab:r3_unitroot}
\begin{threeparttable}
\begin{tabularx}{\linewidth}{>{\raggedright\arraybackslash}Xccc}
\toprule
Variable & LLC statistic & $p$-value & Lags ($\bar{p}$) \\
\midrule
$Y_1$ & $-52.930$ & $<0.001$ & 24 \\
$Y_2$ & $-56.801$ & $<0.001$ & 24 \\
\bottomrule
\end{tabularx}
\begin{tablenotes}[flushleft]
\footnotesize
\item \textit{Notes:} Levin, Lin and Chu (2002) pooled panel unit root test. $H_0$: all panels contain a unit root; $H_1$: all panels are stationary. Model includes individual intercepts (entity demeaning via auxiliary OLS). Lag order selected by AIC over $\bar{p}\in\{0,\ldots,\min(24,\lfloor T^{1/3}\rfloor)\}$ per zone; $\bar{p}$ reports the median across zones. $\bar{T}\approx28,515$ usable observations per zone; at this sample size the LLC finite-sample correction factors (Levin et al.\ 2002, Table~2) are negligible and the statistic is referred to $N(0,1)$ (left-tailed).
\end{tablenotes}
\end{threeparttable}
\end{table}

\end{document}
